\documentclass[11pt]{article}

\usepackage[utf8]{inputenc}
\usepackage{amsmath}
\usepackage{geometry}
\geometry{margin=1in}

\title{Identity Bedrock: A Minimal Substrate for Decentralized Systems}
\author{Jeff Bouchard}
\date{\today}

\begin{document}

\maketitle

\section{Introduction}

Modern networked systems rely on fragmented trust models distributed across identity providers, naming authorities, certificate infrastructures, and centralized service platforms.

\section{Bedrock Identity}

The system is rooted in a minimal identity construct defined by cryptographic key ownership.

\begin{equation}
Identity = (PublicKey, PrivateKey)
\end{equation}

\section{Derived Properties}

All higher-layer functions are modeled as clients of the identity substrate.

\subsection{Naming and Resolution}

Identity enables decentralized naming through a name-value system.

\subsection{Deterministic Execution}

\begin{equation}
Execution = Hash(Input, Environment, Command)
\end{equation}

\section{Conclusion}

A minimal identity substrate can serve as the foundation for a complete network architecture.

\end{document}il 